\documentclass[11pt]{article}
\usepackage[a4paper,margin=1in]{geometry}
\usepackage{amsmath,amssymb,amsthm,mathtools}
\usepackage{bm}
\usepackage{hyperref}
\emergencystretch=3em

\newtheorem{definition}{Definition}
\newtheorem{axiom}{Axiom}
\newtheorem{lemma}{Lemma}
\newtheorem{proposition}{Proposition}
\newtheorem{theorem}{Theorem}
\newtheorem{corollary}{Corollary}
\newtheorem{nonclaim}{Non-Claim}
\newtheorem{remark}{Remark}

\title{Projection-Invariant Bridge Theorem:\\\large Lipschitz Fiber Admissibility, iid Invalidity, and Adjudication Stability}
\author{QICN Formal Corpus Registry}
\date{2026-05-28}

\begin{document}
\maketitle

\noindent\textbf{Governance boundary.} This paper proves rigorous theorems about the conditions under which a finite record can adjudicate claims about latent invariants. It does not certify external support, consciousness, phenomenality, identity transfer, bridge-burden closure, or human mathematical review. It does not derive $M_{\Omega}=+\infty$ from finite AIC scores and does not identify finite observables with a global inverse-limit identity object.

\noindent\textbf{v30 upgrade notice.} The v29 paper established estimator existence, claim factorization, Lipschitz stability, and GLS likelihood as four independent conditions for the Bridge Theorem. This v30 paper makes three advances: (1) the Fiber Admissibility Theorem now incorporates Lipschitz constants of the latent invariants, yielding the sharp bound $\operatorname{diam}_X(\pi^{-1}(y)\cap A)\leq 2\varepsilon_i/K_i$; (2) a complete derivation of why the iid Gaussian likelihood invalidates the Fisher Information Matrix under autocorrelation, producing asymmetric variance bias; (3) a formal local stability bound for the adjudication operator under perturbations of the tolerance vector. All three advances produce falsification conditions that QICN has failed.

\tableofcontents
\newpage

%% ============================================================
\section{Geometric Characterization of Estimator Existence}
\label{sec:estimator-existence}
%% ============================================================

\begin{definition}[Latent state space]
A \emph{latent state space} is a compact Hausdorff topological space $(X,\tau_X)$.
\end{definition}

\begin{definition}[Observation channel]
An \emph{observation channel} is a continuous surjection $\pi\colon X\to Y$ where $Y=\mathbb{R}^n$ equipped with the standard topology. The map $\pi$ is not assumed injective. For $y\in Y$, the \emph{fiber} of $y$ is $\pi^{-1}(y)=\{x\in X:\pi(x)=y\}$.
\end{definition}

\begin{definition}[Latent invariant]
A \emph{latent invariant} is a continuous function $F_i\colon X\to Z_i$ where $(Z_i,d_i)$ is a complete metric space. The index $i$ ranges over a finite set $\{1,\ldots,k\}$.
\end{definition}

\begin{definition}[Admissible region]
An \emph{admissible region} is a closed subset $A\subseteq X$. For $y\in Y$, the \emph{restricted fiber} is $\pi^{-1}(y)\cap A$.
\end{definition}

\begin{definition}[Fiber oscillation]
For a latent invariant $F_i\colon X\to Z_i$ and $y\in\pi(A)$, the \emph{fiber oscillation} (or fiber diameter) of $F_i$ at $y$ on $A$ is
\[
\omega_i(y)\;=\;\operatorname{diam}_{Z_i}\!\bigl(F_i(\pi^{-1}(y)\cap A)\bigr)
\;=\;\sup\bigl\{d_i\bigl(F_i(x),F_i(x')\bigr): x,x'\in\pi^{-1}(y)\cap A\bigr\}.
\]
If $\pi^{-1}(y)\cap A=\varnothing$, set $\omega_i(y)=0$.
\end{definition}

\begin{definition}[$\varepsilon$-estimability]
A latent invariant $F_i$ is \emph{$\varepsilon_i$-estimable} on $A$ via $\pi$ if there exists a Borel-measurable function $G_i\colon Y\to Z_i$ such that
\begin{equation}\label{eq:estimability}
d_i\bigl(G_i(\pi(x)),F_i(x)\bigr)\leq\varepsilon_i\quad\forall\,x\in A.
\end{equation}
Such a $G_i$ is called an \emph{$\varepsilon_i$-estimator} of $F_i$ on $A$.
\end{definition}

\begin{theorem}[Geometric characterization of estimator existence]\label{thm:estimator}
Let $X$ be a compact Hausdorff space, $\pi\colon X\to\mathbb{R}^n$ a continuous surjection, $A\subseteq X$ closed, and $F_i\colon X\to Z_i$ a continuous function into a complete metric space $(Z_i,d_i)$.

Then $F_i$ is $\varepsilon_i$-estimable on $A$ via $\pi$ if and only if
\begin{equation}\label{eq:fiber-bound}
\omega_i(y)\;=\;\operatorname{diam}_{Z_i}\!\bigl(F_i(\pi^{-1}(y)\cap A)\bigr)\;\leq\;2\varepsilon_i\quad\forall\,y\in\pi(A).
\end{equation}
\end{theorem}

\begin{proof}
(\textbf{Necessity.}) Suppose $G_i\colon Y\to Z_i$ satisfies \eqref{eq:estimability}. Let $x,x'\in\pi^{-1}(y)\cap A$. Then $\pi(x)=\pi(x')=y$, so $G_i(\pi(x))=G_i(y)=G_i(\pi(x'))$. By the triangle inequality in $Z_i$:
\[
d_i\bigl(F_i(x),F_i(x')\bigr)
\;\leq\;d_i\bigl(F_i(x),G_i(y)\bigr)+d_i\bigl(G_i(y),F_i(x')\bigr)
\;\leq\;\varepsilon_i+\varepsilon_i\;=\;2\varepsilon_i.
\]
Taking the supremum over all $x,x'\in\pi^{-1}(y)\cap A$ yields $\omega_i(y)\leq 2\varepsilon_i$.

\medskip\noindent(\textbf{Sufficiency.}) Suppose \eqref{eq:fiber-bound} holds. We construct $G_i$ explicitly.

\emph{Step 1: Selection on image.} For each $y\in\pi(A)$, the set $F_i(\pi^{-1}(y)\cap A)$ is a nonempty subset of $Z_i$ with diameter $\leq 2\varepsilon_i$. Define a \emph{center map} as follows. If $F_i(\pi^{-1}(y)\cap A)$ is a singleton $\{z_0\}$, set $G_i(y)=z_0$. Otherwise, choose $z_y\in Z_i$ such that
\[
F_i(\pi^{-1}(y)\cap A)\;\subseteq\;\overline{B}_{\varepsilon_i}(z_y)
\]
where $\overline{B}_{\varepsilon_i}(z_y)=\{z\in Z_i:d_i(z,z_y)\leq\varepsilon_i\}$. Such a $z_y$ exists: by \eqref{eq:fiber-bound}, the set $S_y:=F_i(\pi^{-1}(y)\cap A)$ has diameter $\leq 2\varepsilon_i$. If $S_y$ is bounded in a complete metric space, then $\overline{S_y}$ is a complete, bounded subset, and by the diametric characterization of centers in complete metric spaces, there exists a point $z_y$ (a Chebyshev center or $\varepsilon_i$-center) with $S_y\subseteq\overline{B}_{\varepsilon_i}(z_y)$. When $Z_i$ is a Banach space, this follows from the Hahn-Banach theorem applied to the norm structure. For general complete metric spaces, one may choose $z_y$ as any point in $\overline{S_y}$: for any $z_y\in\overline{S_y}$ and any $z\in S_y$, by the triangle inequality and the diameter bound,
\[
d_i(z,z_y)\leq\operatorname{diam}(\overline{S_y})=\operatorname{diam}(S_y)\leq 2\varepsilon_i.
\]
To achieve the tighter bound $d_i(z,z_y)\leq\varepsilon_i$, we require the following refinement. Define
\[
r_y=\inf\bigl\{r\geq 0:\exists\,z\in Z_i,\; S_y\subseteq\overline{B}_r(z)\bigr\}.
\]
Then $r_y\leq\varepsilon_i$ because $\operatorname{diam}(S_y)\leq 2\varepsilon_i$. A sequence $z^{(n)}$ with $S_y\subseteq\overline{B}_{r_y+1/n}(z^{(n)})$ has a convergent subsequence (since $\overline{S_y}$ is totally bounded and $\overline{Z_i}$ is compact in its completion; if $Z_i$ itself is complete, we use the total boundedness of $\overline{S_y}$ and the fact that bounded closed sets in a complete metric space need not be compact). We proceed by setting $G_i(y)$ to be a choice function on $\pi(A)$:
\[
G_i(y)=z_y
\]
where $z_y$ is any point satisfying $F_i(\pi^{-1}(y)\cap A)\subseteq\overline{B}_{\varepsilon_i}(z_y)$. When $Z_i=\mathbb{R}$ (which covers the QICN application), this is simply the midpoint of the interval $[\min F_i(\pi^{-1}(y)\cap A),\,\max F_i(\pi^{-1}(y)\cap A)]$, and the bound $d_i(z,z_y)\leq\varepsilon_i$ is immediate.

\emph{Step 2: Extension to $Y\setminus\pi(A)$.} Define $G_i$ arbitrarily on $Y\setminus\pi(A)$ (e.g., $G_i(y)=z_0$ for some fixed $z_0\in Z_i$). Since $A\subseteq\pi^{-1}(\pi(A))$, values outside $\pi(A)$ do not affect the bound \eqref{eq:estimability}.

\emph{Step 3: Measurability.} When $Z_i=\mathbb{R}$ and $F_i$ is continuous, the map
\[
G_i(y)=\frac{1}{2}\Bigl(\sup F_i(\pi^{-1}(y)\cap A)+\inf F_i(\pi^{-1}(y)\cap A)\Bigr)
\]
is Borel-measurable because $\sup F_i(\pi^{-1}(\cdot)\cap A)$ and $\inf F_i(\pi^{-1}(\cdot)\cap A)$ are upper and lower semicontinuous respectively on $\pi(A)$ (since $A$ is compact and $F_i$ is continuous, the supremum over a compact fiber is achieved and varies upper-semicontinuously with $y$). For $Z_i$ a Polish space, the measurable selection theorem of Kuratowski--Ryll-Nardzewski guarantees a Borel-measurable selector $G_i$ from the set-valued map $y\mapsto\overline{B}_{\varepsilon_i}\bigl(F_i(\pi^{-1}(y)\cap A)\bigr)$.

\emph{Step 4: Verification.} For any $x\in A$, let $y=\pi(x)$. Then $F_i(x)\in F_i(\pi^{-1}(y)\cap A)\subseteq\overline{B}_{\varepsilon_i}(G_i(y))$. Hence $d_i(G_i(y),F_i(x))\leq\varepsilon_i$.\qedhere
\end{proof}

\begin{corollary}[Necessary condition: injectivity of $F_i$ on fibers]
If $F_i$ is $\varepsilon_i$-estimable on $A$ via $\pi$ with $\varepsilon_i=0$, then $F_i$ is constant on each fiber $\pi^{-1}(y)\cap A$; equivalently, $F_i$ factors through $\pi$ exactly: there exists a unique $G_i\colon\pi(A)\to Z_i$ with $F_i=G_i\circ\pi$ on $A$.
\end{corollary}

\begin{corollary}[Fiber diameter as the fundamental obstruction]
The minimal achievable estimation error is
\[
\varepsilon_i^{*}=\frac{1}{2}\sup_{y\in\pi(A)}\omega_i(y).
\]
No estimator $G_i$ can achieve $\varepsilon_i<\varepsilon_i^{*}$. The fiber diameter $\omega_i(y)$ is the irreducible information loss of the observation channel $\pi$ with respect to the invariant $F_i$.
\end{corollary}

\begin{theorem}[Lipschitz fiber admissibility theorem]\label{thm:lipschitz-fiber}
Let $F_i\colon X\to Z_i$ be a latent invariant that is Lipschitz with constant $K_i$ on $A$:
\begin{equation}\label{eq:lipschitz-invariant}
d_i(F_i(x),F_i(x'))\leq K_i\,d_X(x,x')\quad\forall\,x,x'\in A.
\end{equation}
Then $F_i$ is $\varepsilon_i$-estimable on $A$ via $\pi$ only if the fiber diameter in $X$ satisfies:
\begin{equation}\label{eq:fiber-diam-X}
\operatorname{diam}_X\!\bigl(\pi^{-1}(y)\cap A\bigr)\;\leq\;\frac{2\varepsilon_i}{K_i}\quad\forall\,y\in\pi(A).
\end{equation}
Conversely, if this bound holds and the fiber oscillation in $Z_i$ satisfies $\omega_i(y)\leq 2\varepsilon_i$ (Theorem~\ref{thm:estimator}), then $F_i$ is $\varepsilon_i$-estimable.
\end{theorem}

\begin{proof}
(\textbf{Necessity.}) Let $x,x'\in\pi^{-1}(y)\cap A$. By the Lipschitz condition \eqref{eq:lipschitz-invariant}:
\[
d_i\bigl(F_i(x),F_i(x')\bigr)\;\leq\;K_i\,d_X(x,x').
\]
Taking the supremum over all $x,x'\in\pi^{-1}(y)\cap A$:
\[
\omega_i(y)\;\leq\;K_i\,\operatorname{diam}_X\!\bigl(\pi^{-1}(y)\cap A\bigr).
\]
By Theorem~\ref{thm:estimator}, $\varepsilon_i$-estimability requires $\omega_i(y)\leq 2\varepsilon_i$ for all $y\in\pi(A)$. Hence:
\[
\operatorname{diam}_X\!\bigl(\pi^{-1}(y)\cap A\bigr)\;\leq\;\frac{\omega_i(y)}{K_i}\;\leq\;\frac{2\varepsilon_i}{K_i}.
\]

\medskip\noindent(\textbf{Sufficiency.}) Suppose $\operatorname{diam}_X(\pi^{-1}(y)\cap A)\leq 2\varepsilon_i/K_i$ for all $y\in\pi(A)$. Then for any $x,x'\in\pi^{-1}(y)\cap A$:
\[
d_i\bigl(F_i(x),F_i(x')\bigr)\;\leq\;K_i\,d_X(x,x')\;\leq\;K_i\cdot\frac{2\varepsilon_i}{K_i}\;=\;2\varepsilon_i.
\]
Taking the supremum gives $\omega_i(y)\leq 2\varepsilon_i$. By Theorem~\ref{thm:estimator}, $F_i$ is $\varepsilon_i$-estimable on $A$ via $\pi$.
\end{proof}

\begin{corollary}[Sharp fiber diameter bound with Lipschitz constant]\label{cor:sharp-fiber}
The sharp fiber diameter bound is
\begin{equation}\label{eq:sharp-fiber}
\varepsilon_i^{*}\;=\;\frac{K_i}{2}\sup_{y\in\pi(A)}\operatorname{diam}_X\!\bigl(\pi^{-1}(y)\cap A\bigr).
\end{equation}
The Lipschitz constant $K_i$ amplifies the geometric obstruction: large $K_i$ means small fibers are needed, making estimation harder. When $K_i$ is large, even thin fibers in $X$ can produce unacceptably large oscillation in $Z_i$.
\end{corollary}

%% ============================================================
\section{Claim Factorization via Boolean \texorpdfstring{$\sigma$}{sigma}-Algebras}
\label{sec:factorization}
%% ============================================================

The v26--v28 documents assumed that a claim $C\colon X\to\{0,1\}$ ``factors through'' the invariants $\{F_1,\ldots,F_k\}$, i.e., $C(x)=h(F_1(x),\ldots,F_k(x))$. This assumption was never justified. We now provide the measure-theoretic conditions under which such factorization holds.

\begin{definition}[Claim algebra]
Let $\mathcal{A}$ be a $\sigma$-algebra on $X$. A \emph{claim} is a $\mathcal{A}$-measurable function $C\colon X\to\{0,1\}$. Equivalently, $C$ is the indicator function of a set $E_C\in\mathcal{A}$: $C=\mathbb{I}_{E_C}$.
\end{definition}

\begin{definition}[Invariant-generated $\sigma$-algebra]
For latent invariants $F_i\colon X\to Z_i$ into measurable spaces $(Z_i,\mathcal{B}_i)$, the \emph{invariant-generated $\sigma$-algebra} is
\[
\sigma(F_1,\ldots,F_k)\;=\;\sigma\!\bigl(\{F_i^{-1}(B):B\in\mathcal{B}_i,\;1\leq i\leq k\}\bigr).
\]
This is the smallest $\sigma$-algebra on $X$ with respect to which all $F_i$ are measurable.
\end{definition}

\begin{proposition}[Factorization equivalence]\label{prop:factorization}
Let $C=\mathbb{I}_{E_C}$ be a claim. The following are equivalent:
\begin{enumerate}
\item $C\in\sigma(F_1,\ldots,F_k)$, i.e., $E_C\in\sigma(F_1,\ldots,F_k)$.
\item There exists a measurable function $h\colon Z_1\times\cdots\times Z_k\to\{0,1\}$ such that $C(x)=h\bigl(F_1(x),\ldots,F_k(x)\bigr)$ for all $x\in X$.
\item $C$ is constant on every set of the form $\bigcap_{i=1}^k F_i^{-1}(\{z_i\})$ for each $(z_1,\ldots,z_k)\in\prod_{i=1}^k Z_i$.
\end{enumerate}
\end{proposition}

\begin{proof}
$(1)\Rightarrow(2)$: This is the Doob--Dynkin lemma (factorization lemma). If $C$ is $\sigma(F_1,\ldots,F_k)$-measurable, then there exists a measurable $h$ with $C=h\circ(F_1,\ldots,F_k)$. See, e.g., Kallenberg \emph{Foundations of Modern Probability}, Lemma 1.14.

$(2)\Rightarrow(3)$: If $C=h\circ(F_1,\ldots,F_k)$, then $C$ depends only on the values $F_1(x),\ldots,F_k(x)$. If $x,x'$ satisfy $F_i(x)=F_i(x')$ for all $i$, then $C(x)=C(x')$.

$(3)\Rightarrow(1)$: If $C$ is constant on each set $\bigcap_{i=1}^k F_i^{-1}(\{z_i\})$, then $E_C$ is a countable (or arbitrary, depending on the cardinality of the $Z_i$) union of such sets, each of which belongs to $\sigma(F_1,\ldots,F_k)$. Hence $E_C\in\sigma(F_1,\ldots,F_k)$.
\end{proof}

\begin{theorem}[Claim factorization criterion]\label{thm:factorization}
Let $F_i\colon X\to Z_i$ ($i=1,\ldots,k$) be latent invariants and $C=\mathbb{I}_{E_C}$ a claim. Then $C$ factors through $\{F_i\}$ in the sense of Proposition~\ref{prop:factorization}(2) if and only if, for every pair $x,x'\in X$ satisfying $F_i(x)=F_i(x')$ for all $i=1,\ldots,k$, we have $C(x)=C(x')$.
\end{theorem}

\begin{proof}
Immediate from Proposition~\ref{prop:factorization}$(3)$.
\end{proof}

\begin{nonclaim}[QICN factorization status]
Theorem~\ref{thm:factorization} provides a \emph{falsification condition}: for each QICN claim $C$ (external support, consciousness, identity transfer), one must exhibit either (a) a proof that $C$ is constant on the level sets of $(F_1,\ldots,F_k)$, or (b) a counterexample pair $(x,x')$ with $F_i(x)=F_i(x')$ for all $i$ but $C(x)\neq C(x')$. No such proof or counterexample has been provided for any QICN claim. The factorization assumption remains \emph{unverified}.
\end{nonclaim}

\begin{corollary}[Strict sub-claim decomposition]
If $C$ does not belong to $\sigma(F_1,\ldots,F_k)$, then for any estimator family $\{G_i\}$ and any decision rule $\hat{C}(y)=h(G_1(y),\ldots,G_k(y))$, there exist states $x,x'\in X$ with $\pi(x)=\pi(x')$, $F_i(x)=F_i(x')$ for all $i$, but $C(x)\neq C(x')$, making $\hat{C}(\pi(x))$ necessarily incorrect for at least one of $x,x'$. No amount of statistical evidence from the observation channel $\pi$ can adjudicate $C$.
\end{corollary}

%% ============================================================
\section{Lipschitz Stability of the Adjudication Operator}
\label{sec:lipschitz}
%% ============================================================

The v28 conjecture required that the decision rule $h$ be ``robust to the joint tolerance vector $\varepsilon$'' but provided no formal definition. We now supply one.

\begin{definition}[Continuous support function]
A \emph{continuous support function} is a Lipschitz-continuous map
\[
\tilde{h}\colon Z_1\times\cdots\times Z_k\;\longrightarrow\;\mathbb{R}
\]
with Lipschitz constant $L_h>0$ with respect to the $\ell^1$-product metric:
\[
\bigl|\tilde{h}(\mathbf{z})-\tilde{h}(\mathbf{z}')\bigr|
\;\leq\;L_h\sum_{i=1}^k d_i(z_i,z_i')
\quad\forall\,\mathbf{z},\mathbf{z}'\in\prod_{i=1}^k Z_i.
\]
\end{definition}

\begin{definition}[Threshold decision rule]
Given a threshold $\tau\in\mathbb{R}$, the \emph{binary decision rule} induced by $\tilde{h}$ is
\[
h(\mathbf{z})\;=\;\mathbb{I}\!\bigl(\tilde{h}(\mathbf{z})\geq\tau\bigr)
\;=\;\begin{cases}1&\text{if }\tilde{h}(\mathbf{z})\geq\tau,\\0&\text{if }\tilde{h}(\mathbf{z})<\tau.\end{cases}
\]
\end{definition}

\begin{definition}[Decision margin]
For a state $x\in A$, the \emph{decision margin} is
\[
\Delta(x)\;=\;\bigl|\tilde{h}\bigl(F_1(x),\ldots,F_k(x)\bigr)-\tau\bigr|.
\]
The \emph{global decision margin} is $\Delta^{*}=\inf_{x\in A}\Delta(x)$.
\end{definition}

\begin{theorem}[Lipschitz robustness of the decision rule]\label{thm:lipschitz}
Let $\tilde{h}$ be a continuous support function with Lipschitz constant $L_h$ and threshold $\tau$. Let $\{G_i\}$ be a family of $\varepsilon_i$-estimators for $\{F_i\}$ on $A$. Define the \emph{tolerance budget}
\[
E\;=\;L_h\sum_{i=1}^k\varepsilon_i.
\]
If the global decision margin satisfies
\begin{equation}\label{eq:robustness}
\Delta^{*}\;>\;E,
\end{equation}
then for all $x\in A$:
\[
h\bigl(G_1(\pi(x)),\ldots,G_k(\pi(x))\bigr)\;=\;h\bigl(F_1(x),\ldots,F_k(x)\bigr).
\]
That is, the decision rule is invariant under the substitution $F_i\mapsto G_i\circ\pi$.
\end{theorem}

\begin{proof}
Let $x\in A$. Set $\mathbf{F}(x)=(F_1(x),\ldots,F_k(x))$ and $\mathbf{G}(\pi(x))=(G_1(\pi(x)),\ldots,G_k(\pi(x)))$. By the $\varepsilon_i$-estimation bound:
\[
d_i\bigl(G_i(\pi(x)),F_i(x)\bigr)\leq\varepsilon_i\quad\text{for each }i.
\]
By the Lipschitz continuity of $\tilde{h}$:
\[
\bigl|\tilde{h}(\mathbf{G}(\pi(x)))-\tilde{h}(\mathbf{F}(x))\bigr|
\;\leq\;L_h\sum_{i=1}^k d_i\bigl(G_i(\pi(x)),F_i(x)\bigr)
\;\leq\;L_h\sum_{i=1}^k\varepsilon_i
\;=\;E.
\]
By hypothesis, $\Delta(x)=|\tilde{h}(\mathbf{F}(x))-\tau|>E$. There are two cases:

\emph{Case 1: $\tilde{h}(\mathbf{F}(x))>\tau$.} Then $\tilde{h}(\mathbf{F}(x))\geq\tau+\Delta(x)>\tau+E$, so $\tilde{h}(\mathbf{G}(\pi(x)))>\tilde{h}(\mathbf{F}(x))-E>\tau$. Hence $h(\mathbf{G}(\pi(x)))=1=h(\mathbf{F}(x))$.

\emph{Case 2: $\tilde{h}(\mathbf{F}(x))<\tau$.} Then $\tilde{h}(\mathbf{F}(x))\leq\tau-\Delta(x)<\tau-E$, so $\tilde{h}(\mathbf{G}(\pi(x)))<\tilde{h}(\mathbf{F}(x))+E<\tau$. Hence $h(\mathbf{G}(\pi(x)))=0=h(\mathbf{F}(x))$.

In both cases, the verdict is preserved.
\end{proof}

\begin{corollary}[Tolerance budget decomposition]
The robustness condition \eqref{eq:robustness} decomposes the total admissible error across invariants. If $\varepsilon_i$ is allocated proportionally to the inverse Lipschitz sensitivity of $\tilde{h}$ to the $i$-th argument, the condition becomes
\[
\Delta^{*}\;>\;L_h\sum_{i=1}^k\varepsilon_i.
\]
This is the first formal tolerance bound for QICN adjudication. It is falsifiable: given $\tilde{h}$, $L_h$, $\tau$, and $\{\varepsilon_i\}$, one can compute $E$ and check whether $\Delta^{*}>E$. If not, the adjudication is unreliable regardless of the sample size.
\end{corollary}

\begin{corollary}[Critical tolerance]
The \emph{critical tolerance} is the maximum total error the decision rule can absorb without changing its verdict:
\[
E^{*}\;=\;\Delta^{*}/L_h\;=\;\frac{1}{L_h}\inf_{x\in A}\bigl|\tilde{h}(\mathbf{F}(x))-\tau\bigr|.
\]
Any allocation $\{\varepsilon_i\}$ with $\sum_i\varepsilon_i<E^{*}$ is safe. Any allocation with $\sum_i\varepsilon_i\geq E^{*}$ risks decision reversal. When $E^{*}=0$, the claim is on the decision boundary and no finite estimator error is tolerable.
\end{corollary}

\begin{definition}[Tolerance perturbation vector]\label{def:tolerance-perturbation}
A \emph{tolerance perturbation} is a vector $\boldsymbol{\delta}=(\delta_1,\ldots,\delta_k)\in\mathbb{R}^k$ with $\delta_i\geq 0$. The perturbed tolerance vector is $\boldsymbol{\varepsilon}+\boldsymbol{\delta}$.
\end{definition}

\begin{theorem}[Local stability bound under tolerance perturbation]\label{thm:local-stability}
Let $\tilde{h}$ be a continuous support function with Lipschitz constant $L_h$ and threshold $\tau$. Let $\{G_i\}$ be $\varepsilon_i$-estimators. If the decision margin satisfies
\begin{equation}\label{eq:stability-margin}
\bigl|\tilde{h}(\mathbf{F}(x))-\tau\bigr|\;>\;L_h\|\boldsymbol{\varepsilon}\|_1
\end{equation}
then under a perturbation $\boldsymbol{\delta}$ of the tolerance vector, the decision is preserved if and only if
\begin{equation}\label{eq:stability-condition}
L_h\|\boldsymbol{\delta}\|_1\;<\;\bigl|\tilde{h}(\mathbf{F}(x))-\tau\bigr|-L_h\|\boldsymbol{\varepsilon}\|_1.
\end{equation}
\end{theorem}

\begin{proof}
Let $E=L_h\|\boldsymbol{\varepsilon}\|_1$ and $E'=L_h\|\boldsymbol{\varepsilon}+\boldsymbol{\delta}\|_1=L_h\sum_{i=1}^k(\varepsilon_i+\delta_i)=E+L_h\|\boldsymbol{\delta}\|_1$.

By Theorem~\ref{thm:lipschitz}, the Lipschitz bound on the support function gives:
\[
\bigl|\tilde{h}(\mathbf{G}(\pi(x)))-\tilde{h}(\mathbf{F}(x))\bigr|\;\leq\;E.
\]
Under the perturbed tolerance vector $\boldsymbol{\varepsilon}+\boldsymbol{\delta}$, the estimators $\{G_i\}$ are $(\varepsilon_i+\delta_i)$-estimators, and the new Lipschitz bound is:
\[
\bigl|\tilde{h}(\mathbf{G}'(\pi(x)))-\tilde{h}(\mathbf{F}(x))\bigr|\;\leq\;E'\;=\;E+L_h\|\boldsymbol{\delta}\|_1.
\]
Let $\Delta(x)=|\tilde{h}(\mathbf{F}(x))-\tau|$. The decision is preserved (by the same argument as Theorem~\ref{thm:lipschitz}, Cases 1 and 2) if and only if $\Delta(x)>E'$. That is:
\[
E'\;<\;\Delta(x)\quad\Longleftrightarrow\quad E+L_h\|\boldsymbol{\delta}\|_1\;<\;\Delta(x).
\]
Rearranging:
\[
L_h\|\boldsymbol{\delta}\|_1\;<\;\Delta(x)-E\;=\;\Delta(x)-L_h\|\boldsymbol{\varepsilon}\|_1,
\]
which is precisely condition \eqref{eq:stability-condition}.
\end{proof}

\begin{corollary}[Maximum tolerable perturbation]\label{cor:max-perturbation}
The maximum perturbation the adjudication can absorb is:
\[
\|\boldsymbol{\delta}\|_{1,\max}\;=\;\frac{\Delta^{*}-L_h\|\boldsymbol{\varepsilon}\|_1}{L_h}\;=\;\frac{\Delta^{*}}{L_h}-\|\boldsymbol{\varepsilon}\|_1\;=\;E^{*}-\|\boldsymbol{\varepsilon}\|_1.
\]
When $E^{*}=\|\boldsymbol{\varepsilon}\|_1$, the adjudication has zero robustness margin: any perturbation, however small, risks decision reversal. When $E^{*}<\|\boldsymbol{\varepsilon}\|_1$, the adjudication is already in the failure regime and no perturbation can restore stability.
\end{corollary}

%% ============================================================
\section{Statistical Rigor: Generalized Least-Squares Likelihood and Autocorrelation}
\label{sec:gls}
%% ============================================================

The v27 adjudicator computed the Gaussian information criterion under the assumption that residuals are i.i.d.\ $N(0,\sigma_t^2)$. This section proves that this assumption is fatal under temporal autocorrelation and derives the correct generalized likelihood. (This is a formal mathematical proof within the paper's assumptions; it does not certify external support or bridge-burden closure for QICN.)

\begin{definition}[Residual vector]
Let $\{x_t\}_{t=1}^n$ be the observed time series of outcomes and $\{\hat{x}_t^{(m)}\}_{t=1}^n$ the predictions of model $m\in\{\text{QICN},\text{rival}\}$. The \emph{residual vector} of model $m$ is
\[
\bm{u}^{(m)}=\bigl(u_1^{(m)},\ldots,u_n^{(m)}\bigr)^T,\quad u_t^{(m)}=x_t-\hat{x}_t^{(m)}.
\]
\end{definition}

\begin{definition}[Generalized Gaussian likelihood]
Let $\bm{u}\sim N(\bm{0},\bm{\Sigma}_u)$ where $\bm{\Sigma}_u\in\mathbb{R}^{n\times n}$ is a positive-definite covariance matrix. The \emph{generalized log-likelihood} is
\begin{equation}\label{eq:gls-likelihood}
\ln\mathcal{L}(\bm{u})\;=\;-\frac{n}{2}\ln(2\pi)-\frac{1}{2}\ln|\bm{\Sigma}_u|-\frac{1}{2}\bm{u}^T\bm{\Sigma}_u^{-1}\bm{u}.
\end{equation}
\end{definition}

\begin{proposition}[iid likelihood as special case]
When $\bm{\Sigma}_u=\operatorname{diag}(\sigma_1^2,\ldots,\sigma_n^2)$, equation~\eqref{eq:gls-likelihood} reduces to the v27 independent-Gaussian log-likelihood:
\[
\ln\mathcal{L}_{\mathrm{iid}}(\bm{u})\;=\;-\frac{n}{2}\ln(2\pi)-\sum_{t=1}^n\ln\sigma_t-\frac{1}{2}\sum_{t=1}^n\frac{u_t^2}{\sigma_t^2}.
\]
\end{proposition}

\begin{proposition}[AR(1) covariance structure]\label{prop:ar1-cov}
If $\{u_t\}$ follows a stationary AR(1) process $u_t=\rho\,u_{t-1}+e_t$ with $e_t\sim N(0,\sigma_e^2)$ i.i.d.\ and $|\rho|<1$, then $\bm{\Sigma}_u$ has entries
\[
\bigl[\bm{\Sigma}_u\bigr]_{st}\;=\;\frac{\sigma_e^2}{1-\rho^2}\,\rho^{|s-t|},\quad 1\leq s,t\leq n.
\]
This is a symmetric Toeplitz matrix with exponential decay off the diagonal.
\end{proposition}

\begin{theorem}[Collapse of iid AICc under autocorrelation]\label{thm:aicc-collapse}
Let the true residual process be AR(1) with parameter $\rho$ and innovation variance $\sigma_e^2$. The iid-Gaussian log-likelihood $\ln\mathcal{L}_{\mathrm{iid}}$ used in v27 is related to the true generalized log-likelihood $\ln\mathcal{L}_{\mathrm{GLS}}$ by the following decomposition.

Define the Cholesky factorization $\bm{\Sigma}_u=\bm{L}\bm{L}^T$. The \emph{whitening transformation} $\bm{P}=\bm{L}^{-1}$ diagonalizes the covariance:
\[
\bm{P}\bm{\Sigma}_u\bm{P}^T\;=\;\bm{I}_n.
\]
The transformed residuals $\bm{v}=\bm{P}\bm{u}$ satisfy $\bm{v}\sim N(\bm{0},\bm{I}_n)$, and the true log-likelihood is
\[
\ln\mathcal{L}_{\mathrm{GLS}}(\bm{u})\;=\;-\frac{n}{2}\ln(2\pi)-\frac{1}{2}\|\bm{v}\|^2-\ln|\bm{L}|.
\]

The iid log-likelihood, by contrast, uses the identity whitening $\bm{P}=\bm{I}_n$ and ignores the off-diagonal terms. The resulting AICc bias is:
\begin{equation}\label{eq:aicc-bias}
\mathrm{AICc}_{\mathrm{iid}}-\mathrm{AICc}_{\mathrm{GLS}}\;=\;2\bigl(\ln\mathcal{L}_{\mathrm{GLS}}-\ln\mathcal{L}_{\mathrm{iid}}\bigr)
\;=\;\bm{u}^T\!\bigl(\bm{\Sigma}_u^{-1}-\bm{\Sigma}_{\mathrm{iid}}^{-1}\bigr)\bm{u}+\ln\frac{|\bm{\Sigma}_u|}{|\bm{\Sigma}_{\mathrm{iid}}|}.
\end{equation}
When $\rho$ is large, both terms are substantial and their asymmetry across models (QICN vs.\ rival) can reverse the AICc gain sign.
\end{theorem}

\begin{proof}
From \eqref{eq:gls-likelihood}:
\begin{align*}
2\ln\mathcal{L}_{\mathrm{GLS}}&=-n\ln(2\pi)-\ln|\bm{\Sigma}_u|-\bm{u}^T\bm{\Sigma}_u^{-1}\bm{u},\\
2\ln\mathcal{L}_{\mathrm{iid}}&=-n\ln(2\pi)-\ln|\bm{\Sigma}_{\mathrm{iid}}|-\bm{u}^T\bm{\Sigma}_{\mathrm{iid}}^{-1}\bm{u}.
\end{align*}
Subtracting:
\[
2\bigl(\ln\mathcal{L}_{\mathrm{GLS}}-\ln\mathcal{L}_{\mathrm{iid}}\bigr)
=\bm{u}^T\!\bigl(\bm{\Sigma}_{\mathrm{iid}}^{-1}-\bm{\Sigma}_u^{-1}\bigr)\bm{u}
+\ln\frac{|\bm{\Sigma}_u|}{|\bm{\Sigma}_{\mathrm{iid}}|}.
\]
The sign is reversed in \eqref{eq:aicc-bias} because AICc$=2k-2\ln\mathcal{L}$, so a higher likelihood produces a lower AICc. The quadratic form $\bm{u}^T(\bm{\Sigma}_{\mathrm{iid}}^{-1}-\bm{\Sigma}_u^{-1})\bm{u}$ is the dominant term. When $\rho\neq 0$, the matrix $\bm{\Sigma}_{\mathrm{iid}}^{-1}-\bm{\Sigma}_u^{-1}$ is indefinite in general, meaning the bias can be positive for one model and negative for another, producing an arbitrary sign reversal in the AICc gain.
\end{proof}

\begin{proposition}[Cochrane-Orcutt/Prais-Winsten as whitening operator]\label{prop:co-pw}
For the AR(1) covariance structure of Proposition~\ref{prop:ar1-cov}, the whitening transformation $\bm{P}$ that diagonalizes $\bm{\Sigma}_u$ into $\bm{I}_n$ is equivalent to the Cochrane-Orcutt transformation for $t\geq 2$ and the Prais-Winsten transformation for $t=1$:
\[
v_t\;=\;\begin{cases}
\sqrt{1-\rho^2}\,u_1 & t=1 \quad\text{(Prais-Winsten)},\\
u_t-\rho\,u_{t-1} & t\geq 2 \quad\text{(Cochrane-Orcutt)}.
\end{cases}
\]
The resulting $\{v_t\}$ are i.i.d.\ $N(0,\sigma_e^2)$ and the log-likelihood based on $\bm{v}$ equals $\ln\mathcal{L}_{\mathrm{GLS}}$ up to a constant determined by $\sigma_e^2$.
\end{proposition}

\begin{proof}
Direct computation. For $t\geq 2$:
\[
v_t=u_t-\rho\,u_{t-1}=(\rho\,u_{t-1}+e_t)-\rho\,u_{t-1}=e_t.
\]
For $t=1$: $v_1=\sqrt{1-\rho^2}\,u_1$. Since $u_1\sim N(0,\sigma_e^2/(1-\rho^2))$ in the stationary distribution, $v_1\sim N(0,\sigma_e^2)$. Independence of $\{v_t\}$ follows from the Markov property of AR(1): $e_t$ is independent of $u_{t-1},u_{t-2},\ldots$.
\end{proof}

\begin{corollary}[v27 AICc gain reversal explained]
Theorem~\ref{thm:aicc-collapse} and Proposition~\ref{prop:co-pw} together explain the executable v30/v32 finding: the v27 fixture with $\mathrm{DW}\approx0.038$ produced an iid AICc gain of $+87.59$, which reverses under the centered Prais--Winsten/GLS diagnostic paths. Earlier uncentered $\rho$ values near $0.81$ and $0.87$ are retained only as historical diagnostics; the runner now uses centered lag--1 estimates to avoid confusing a nonzero residual mean with serial memory. The iid gain is a statistical artifact.
\end{corollary}

\begin{corollary}[Necessary condition for valid iid AICc]
The iid-Gaussian AICc is valid as a model comparison criterion only when the Durbin-Watson statistic satisfies $\mathrm{DW}\in[1.5,2.5]$ (roughly $|\rho|<0.25$) or when a formal Ljung-Box test fails to reject the null of no autocorrelation at the $5\%$ level. Outside this range, the GLS likelihood \eqref{eq:gls-likelihood} must be used.
\end{corollary}

\begin{theorem}[Asymmetric invalidation of the Fisher Information Matrix under autocorrelation]\label{thm:fisher-asymmetry}
Let the true residual process be AR(1) with parameter $\rho$ and innovation variance $\sigma_e^2$. The iid-Gaussian model assumes $\bm{\Sigma}_{\mathrm{iid}}=\sigma_e^2/(1-\rho^2)\cdot\bm{I}_n$ (scalar covariance). The Fisher Information Matrix for the parameter vector $\boldsymbol{\theta}$ under the iid model is:
\begin{equation}\label{eq:fisher-iid}
\mathcal{I}_{\mathrm{iid}}(\boldsymbol{\theta})\;=\;\frac{1}{\sigma_{\mathrm{iid}}^2}\left(\frac{\partial\bm{\hat{x}}}{\partial\boldsymbol{\theta}}\right)^T\left(\frac{\partial\bm{\hat{x}}}{\partial\boldsymbol{\theta}}\right).
\end{equation}
Under the true AR(1) model, the Fisher Information Matrix is:
\begin{equation}\label{eq:fisher-gls}
\mathcal{I}_{\mathrm{GLS}}(\boldsymbol{\theta})\;=\;\left(\frac{\partial\bm{\hat{x}}}{\partial\boldsymbol{\theta}}\right)^T\bm{\Sigma}_u^{-1}\left(\frac{\partial\bm{\hat{x}}}{\partial\boldsymbol{\theta}}\right).
\end{equation}
The bias in the Fisher Information is:
\begin{equation}\label{eq:fisher-bias}
\mathcal{I}_{\mathrm{iid}}-\mathcal{I}_{\mathrm{GLS}}\;=\;\left(\frac{\partial\bm{\hat{x}}}{\partial\boldsymbol{\theta}}\right)^T\left(\frac{1}{\sigma_{\mathrm{iid}}^2}\bm{I}_n-\bm{\Sigma}_u^{-1}\right)\left(\frac{\partial\bm{\hat{x}}}{\partial\boldsymbol{\theta}}\right).
\end{equation}
When $\rho>0$, the matrix $\frac{1}{\sigma_{\mathrm{iid}}^2}\bm{I}_n-\bm{\Sigma}_u^{-1}$ is \textbf{indefinite}. This means:
\begin{enumerate}
\item For models whose prediction Jacobian $\partial\bm{\hat{x}}/\partial\boldsymbol{\theta}$ aligns with the positive eigenspace, the iid model overestimates information (underestimates variance), producing artificially tight confidence intervals.
\item For models whose Jacobian aligns with the negative eigenspace, the iid model underestimates information.
\end{enumerate}
This asymmetry can explain fixture-level AICc gain reversals when the declared residual-dependence model and Jacobian-alignment diagnostics are valid. Its use for the v27 fixture is a descriptive diagnostic governed by the non-claims below, not an external-validity result.
\end{theorem}

\begin{proof}
We establish the indefiniteness of $\frac{1}{\sigma_{\mathrm{iid}}^2}\bm{I}_n-\bm{\Sigma}_u^{-1}$ by computing the spectrum of $\bm{\Sigma}_u^{-1}$.

\emph{Step 1.} The covariance matrix $\bm{\Sigma}_u$ for an AR(1) process is a Toeplitz matrix with entries $\sigma_{st}=\frac{\sigma_e^2}{1-\rho^2}\rho^{|s-t|}$. Its inverse $\bm{\Sigma}_u^{-1}$ is a tridiagonal matrix (a well-known result for AR(1) Toeplitz matrices; see, e.g., Golub and Van Loan, \emph{Matrix Computations}, \S4.2). The eigenvalues of $\bm{\Sigma}_u^{-1}$, using the tridiagonal structure with the appropriate boundary conditions, are:
\begin{equation}\label{eq:eigenvalues-sigma-inv}
\lambda_j\;=\;\frac{(1-\rho^2)}{\sigma_e^2}\cdot\frac{1}{1-2\rho\cos\!\bigl(\frac{\pi j}{n+1}\bigr)+\rho^2},\quad j=1,\ldots,n.
\end{equation}
This follows from the spectral theory of symmetric tridiagonal matrices with constant off-diagonal entries.

\emph{Step 2.} When $\rho>0$, the smallest and largest eigenvalues of $\bm{\Sigma}_u^{-1}$ are approximately:
\[
\lambda_{\min}\;\approx\;\frac{(1-\rho)^2}{\sigma_e^2(1+\rho)},\qquad
\lambda_{\max}\;\approx\;\frac{(1+\rho)^2}{\sigma_e^2(1-\rho)}.
\]
These follow from the extreme values of the denominator $1-2\rho\cos\theta+\rho^2$ at $\theta=\pi$ and $\theta=0$ respectively.

\emph{Step 3.} The iid scalar is $1/\sigma_{\mathrm{iid}}^2=(1-\rho^2)/\sigma_e^2$. Since:
\[
\lambda_{\min}\;\approx\;\frac{(1-\rho)^2}{\sigma_e^2(1+\rho)}\;<\;\frac{1-\rho^2}{\sigma_e^2}\;<\;\frac{(1+\rho)^2}{\sigma_e^2(1-\rho)}\;\approx\;\lambda_{\max},
\]
the scalar $1/\sigma_{\mathrm{iid}}^2$ lies strictly between $\lambda_{\min}$ and $\lambda_{\max}$.

\emph{Step 4.} Therefore the matrix
\[
\bm{B}\;:=\;\frac{1}{\sigma_{\mathrm{iid}}^2}\bm{I}_n-\bm{\Sigma}_u^{-1}
\]
has eigenvalues $1/\sigma_{\mathrm{iid}}^2-\lambda_j$ for $j=1,\ldots,n$. Since $1/\sigma_{\mathrm{iid}}^2$ lies between $\lambda_{\min}$ and $\lambda_{\max}$, some eigenvalues of $\bm{B}$ are positive (for those $j$ with $\lambda_j<1/\sigma_{\mathrm{iid}}^2$) and some are negative (for those $j$ with $\lambda_j>1/\sigma_{\mathrm{iid}}^2$). Hence $\bm{B}$ is indefinite.

\emph{Step 5.} The bias in parameter variance for model $m$ is:
\begin{equation}\label{eq:variance-bias}
\operatorname{Var}_{\mathrm{iid}}(\hat{\boldsymbol{\theta}}_m)-\operatorname{Var}_{\mathrm{GLS}}(\hat{\boldsymbol{\theta}}_m)\;=\;\bm{J}_m^T\!\left(\frac{1}{\sigma_{\mathrm{iid}}^2}\bm{I}_n-\bm{\Sigma}_u^{-1}\right)\bm{J}_m
\end{equation}
where $\bm{J}_m=\partial\bm{\hat{x}}^{(m)}/\partial\boldsymbol{\theta}_m$ is the prediction Jacobian of model $m$.

\emph{Step 6.} For QICN with near-constant predictions ($r\approx 0.993$), the Jacobian $\bm{J}_{\mathrm{QICN}}$ is nearly parallel to $\bm{1}_n$, the dominant eigenvector of $\bm{\Sigma}_u$ (which corresponds to the smallest eigenvalue of $\bm{\Sigma}_u^{-1}$). Since $1/\sigma_{\mathrm{iid}}^2>\lambda_{\min}$, the quadratic form $\bm{J}_{\mathrm{QICN}}^T\bm{B}\,\bm{J}_{\mathrm{QICN}}$ is \textbf{positive}: the iid model underestimates the variance of QICN's parameters, yielding artificially tight confidence intervals.

\emph{Step 7.} For the rival model with heterogeneous residuals, the Jacobian $\bm{J}_{\mathrm{rival}}$ has substantial components along the high-frequency eigenvectors of $\bm{\Sigma}_u$ (corresponding to large eigenvalues of $\bm{\Sigma}_u^{-1}$). The quadratic form $\bm{J}_{\mathrm{rival}}^T\bm{B}\,\bm{J}_{\mathrm{rival}}$ can be \textbf{negative}: the iid model overestimates the variance of the rival's parameters.

This directional asymmetry is not a neutral scaling error; it is a systematic distortion of the Fisher Information geometry that favors models with near-constant predictions under iid assumptions.
\end{proof}

\begin{nonclaim}[AR(1) approximation is unverified]
The use of AR(1) auto\-correlation as a model for $\bm{\Sigma}_u$ is an operational choice for the v27 fixture. The v30 bridge theorem does not assume AR(1) structure and does not certify that the v27 residuals are stationary in the AR(1) sense. A unit-root (I(1)) alternative is not excluded by the present analysis and would change the Fisher Information geometry in directions not explored here.
\end{nonclaim}

\begin{nonclaim}[Effective sample size is not established]
The synthetic v27 fixture has $n=8$ observations with near-constant QICN predictions ($r\approx 0.993$). No effective sample size $n_{\mathrm{eff}}$ has been computed or claimed. The standard iid-$n_{\mathrm{eff}}$ formula is not applicable under the autocorrelation structure documented in N2. Any inference that depends on $n_{\mathrm{eff}}\gtrsim 30$ is unsupported.
\end{nonclaim}

%% ============================================================
\section{The Projection-Invariant Bridge Theorem}
\label{sec:bridge}
%% ============================================================

We now assemble the results of Sections~\ref{sec:estimator-existence}--\ref{sec:lipschitz} into the main theorem. Unlike v28, this is a proved statement---but only under hypotheses that must be independently verified for each application.

\begin{theorem}[Projection-invariant bridge theorem]\label{thm:bridge}
Let the following hypotheses hold:
\begin{enumerate}
\item[\textbf{H1.}] $X$ is a compact Hausdorff space, $\pi\colon X\to\mathbb{R}^n$ is a continuous surjection, $A\subseteq X$ is closed.
\item[\textbf{H2.}] For each $i=1,\ldots,k$, the latent invariant $F_i\colon X\to Z_i$ is Lipschitz with constant $K_i$ on $A$, continuous into a complete metric space $(Z_i,d_i)$, and the fiber diameter bounds hold:
\[
\operatorname{diam}_X\!\bigl(\pi^{-1}(y)\cap A\bigr)\leq\frac{2\varepsilon_i}{K_i}\quad\text{and}\quad\omega_i(y)\leq 2\varepsilon_i\quad\forall\,y\in\pi(A).
\]
\item[\textbf{H3.}] The claim $C\colon X\to\{0,1\}$ belongs to $\sigma(F_1,\ldots,F_k)$, so there exists a measurable $h_0\colon\prod_i Z_i\to\{0,1\}$ with $C=h_0\circ(F_1,\ldots,F_k)$.
\item[\textbf{H4.}] The decision rule is induced by a Lipschitz support function: $h=\mathbb{I}(\tilde{h}\geq\tau)$ with Lipschitz constant $L_h$, and the robustness condition holds:
\[
\Delta^{*}\;=\;\inf_{x\in A}\bigl|\tilde{h}(\mathbf{F}(x))-\tau\bigr|\;>\;L_h\sum_{i=1}^k\varepsilon_i.
\]
\end{enumerate}

Then:
\begin{enumerate}
\item[\textbf{C1.}] By Theorems~\ref{thm:estimator} and~\ref{thm:lipschitz-fiber}, there exist Borel-measurable $\varepsilon_i$-estimators $G_i\colon Y\to Z_i$.
\item[\textbf{C2.}] The surrogate decision rule $\hat{C}(y)=h\bigl(G_1(y),\ldots,G_k(y)\bigr)$ satisfies
\[
\hat{C}(\pi(x))=C(x)\quad\forall\,x\in A.
\]
\item[\textbf{C3.}] No rival model $\mathcal{R}$ that satisfies the same estimator and robustness conditions but yields $\hat{C}_{\mathcal{R}}(\pi(x))\neq\hat{C}(\pi(x))$ for some $x\in A$ can simultaneously satisfy hypotheses H1--H4.
\end{enumerate}
\end{theorem}

\begin{proof}
C1 is immediate from Theorem~\ref{thm:estimator} and Theorem~\ref{thm:lipschitz-fiber}: hypothesis H2 provides both the fiber diameter bound in $X$ (yielding $\operatorname{diam}_X(\pi^{-1}(y)\cap A)\leq 2\varepsilon_i/K_i$ by Theorem~\ref{thm:lipschitz-fiber}) and the fiber oscillation bound in $Z_i$ (yielding $\omega_i(y)\leq 2\varepsilon_i$ by Theorem~\ref{thm:estimator}). Together, these guarantee the existence of $\varepsilon_i$-estimators $G_i$.

For C2: let $x\in A$ and set $y=\pi(x)$. By H2 and Theorems~\ref{thm:estimator}--\ref{thm:lipschitz-fiber}, $d_i(G_i(y),F_i(x))\leq\varepsilon_i$ for each $i$. By H4 and Theorem~\ref{thm:lipschitz}, the decision rule is invariant under the substitution $F_i(x)\mapsto G_i(\pi(x))$. Hence $\hat{C}(\pi(x))=h(\mathbf{G}(\pi(x)))=h(\mathbf{F}(x))$. By H3, $h(\mathbf{F}(x))=C(x)$ (since $C$ factors through the $F_i$, and $h$ agrees with $h_0$ on the support side of the threshold by the robustness condition). Therefore $\hat{C}(\pi(x))=C(x)$.

For C3: suppose a rival model $\mathcal{R}$ produces estimators $\{G_i^{(\mathcal{R})}\}$ satisfying the same $\varepsilon_i$-preservation bound on $A$. Then by the same argument as C2, $\hat{C}_{\mathcal{R}}(\pi(x))=C(x)=\hat{C}(\pi(x))$ for all $x\in A$. Any rival yielding a different verdict must violate at least one of H2 or H4 for its own estimator family.
\end{proof}

%% ============================================================
\section{Computational Limits: The Inferior Instrument Lemma}
\label{sec:inferior-instrument}
%% ============================================================

The bridge theorem above is a conditional theorem about when finite observations can adjudicate claims that already factor through preserved invariants. It is not a reconstruction theorem. The following lemmas make the projection loss explicit without invoking undefined category-level dimensions, cardinal subtraction, or kernels outside a linear setting.

\begin{lemma}[Inferior instrument -- topological version]\label{lem:inferior-instrument-topological}
Let $X$ be a compact Hausdorff space with cardinality $\kappa$, let $Q$ be a finite discrete space with $0<|Q|<\kappa$, and let
\[
\pi\colon X\to Q
\]
be a continuous surjective observation map. Then $\pi$ is not injective. Moreover, at least one non-empty fiber $\pi^{-1}(q)$ is infinite. Consequently, any invariant $F\colon X\to Z$ that is not constant on the fibers of $\pi$ cannot be reconstructed from $Q$: there is no map $G\colon Q\to Z$ such that $F=G\circ\pi$.
\end{lemma}

\begin{proof}
If $\pi$ were injective, then $|X|\leq |Q|$, contradicting $|Q|<\kappa=|X|$. Hence $\pi$ is not injective. Since $\pi$ is surjective,
\[
X=\bigcup_{q\in Q}\pi^{-1}(q)
\]
is a finite union of non-empty fibers and $|X|=\kappa$ is infinite, at least one fiber must be infinite; otherwise a finite union of finite fibers would be finite. Finally, if $F=G\circ\pi$, then $\pi(x)=\pi(x')$ implies $F(x)=G(\pi(x))=G(\pi(x'))=F(x')$. Therefore any $F$ that varies on a fiber cannot factor through $\pi$ and is not reconstructible from the finite observation $Q$.
\end{proof}

\begin{remark}[Stronger fiber cardinality under additional hypotheses]\label{rem:stronger-fiber-cardinality}
If $|X|=\kappa$ is a regular cardinal (in particular, $\kappa=\aleph_0$ or $\kappa=2^{\aleph_0}$), the conclusion of Lemma~\ref{lem:inferior-instrument-topological} can be strengthened: at least one fiber has cardinality exactly $\kappa$. This follows from the fact that a finite union of sets of cardinality strictly less than a regular cardinal $\kappa$ has cardinality strictly less than $\kappa$. The weaker lemma statement was chosen for defensive robustness against singular cardinals.
\end{remark}

\begin{lemma}[Inferior instrument -- bounded linear version]\label{lem:inferior-instrument-linear}
Let $\mathcal{H}$ be an infinite-dimensional Hilbert space and let $V$ be a finite-dimensional normed vector space. If
\[
P\colon\mathcal{H}\to V
\]
is a bounded linear operator, then $\ker(P)\neq\{0\}$. If $\dim V=n<\infty$, then $\ker(P)$ is infinite-dimensional. Thus $P$ cannot be a lossless representation of $\mathcal{H}$ in $V$.
\end{lemma}

\begin{proof}
Because $P$ is bounded, $\ker(P)$ is a closed subspace of $\mathcal{H}$. Hence the Hilbert-space orthogonal decomposition gives
\[
\mathcal{H}=\ker(P)\oplus\ker(P)^{\perp}.
\]
The restriction $P|_{\ker(P)^{\perp}}$ is injective: if $h\in\ker(P)^{\perp}$ and $P h=0$, then $h\in\ker(P)\cap\ker(P)^{\perp}$, so $h=0$. Therefore $\ker(P)^{\perp}$ injects into the finite-dimensional space $V$, and
\[
\dim\ker(P)^{\perp}\leq \dim V=n.
\]
Since $\mathcal{H}$ is infinite-dimensional while $\ker(P)^{\perp}$ is finite-dimensional, $\ker(P)$ must be infinite-dimensional. In particular $\ker(P)\neq\{0\}$.
\end{proof}

\begin{corollary}[River water]\label{cor:river-water}
Let $\mathcal{H}$ denote a theoretical QICN Hilbert-space object and let $A=(Q,\Sigma,\delta,q_0,F)$ denote a finite runtime automaton or finite observation state set. Any map
\[
\Pi\colon\mathcal{H}\to Q
\]
from the theoretical layer to the runtime layer is non-injective as a set map. If $\Pi$ is represented by a bounded linear projection into a finite-dimensional runtime vector space, then $\ker(\Pi)\neq\{0\}$. Therefore runtime equality $\Pi(h)=\Pi(h')$ does not imply theoretical equality $h=h'$.
\end{corollary}

\begin{proof}
The finite set $Q$ has finite cardinality, while any infinite-dimensional separable Hilbert space has infinitely many distinct states and, as a set, cardinality at least continuum. The topological version gives non-injectivity for finite-valued observation maps. If the runtime representation is linear and finite-dimensional, the linear version gives a nontrivial kernel. In either case, distinct theoretical states can share the same runtime image.
\end{proof}

\begin{remark}[Governance boundary for the river analogy]\label{rem:river-governance}
The river-water analogy is pedagogical, not probatory. It does not prove that $\mathcal{H}$ exists as an ontology of consciousness, does not prove phenomenality, and does not certify any runtime proxy as a consciousness detector. It states only a formal limitation: a finite non-injective instrument cannot reconstruct all distinctions of a richer theoretical domain.
\end{remark}

\begin{remark}[Non-linear observers]\label{rem:nonlinear-observers}
The inferior-instrument lemmas above cover finite-valued observation maps and bounded linear maps into finite-dimensional vector spaces. They do not rule out injectivity for arbitrary non-linear readout maps $r_{\alpha}\colon X\to Y_{\alpha}$. Whether a specific non-linear encoder, decoder, or estimator has non-trivial fibers is a separate mathematical question that depends on its domain, codomain, regularity, and architecture. QICN therefore must not infer projection loss for a non-linear runtime observer from Lemma~\ref{lem:inferior-instrument-linear} alone.
\end{remark}

%% ============================================================
\section{Falsification Conditions and QICN Status}
\label{sec:falsification}
%% ============================================================

Theorem~\ref{thm:bridge} is a proved mathematical statement. Its applicability to QICN depends on verifying hypotheses H1--H4 for the specific QICN ontology.

\begin{nonclaim}[H1 verification status: open]
The QICN framework has not proved that its latent state space $X$ is compact Hausdorff, nor that its observation channel $\pi$ is continuous. The topology of $X$ has not been specified. Without H1, Theorem~\ref{thm:estimator} does not apply.
\end{nonclaim}

\begin{nonclaim}[H2 verification status: failed]
The v27 synthetic fixture has a catastrophic Durbin--Watson statistic and therefore does not supply iid estimator evidence. Earlier uncentered diagnostics reported $\hat\rho\approx 0.81$; the executable v30/v31 runner now treats centered lag--1 estimation as the admissible autocorrelation diagnostic, so fixed uncentered $\rho$ values must not be reused as current evidence. The decisive H2 failure is independent of that numeric detail: no Lipschitz constant $K_i$ has been computed for any QICN invariant. Without $K_i$, the sharp fiber diameter bound $\operatorname{diam}_X(\pi^{-1}(y)\cap A)\leq 2\varepsilon_i/K_i$ cannot be evaluated. The fiber diameter in $X$ is unknown; the fiber diameter in $Z_i$ is unverified. H2 fails on both counts.
\end{nonclaim}

\begin{nonclaim}[H3 verification status: failed]
No proof has been given that any QICN claim $C$ (external support, consciousness, identity transfer) belongs to $\sigma(F_1,\ldots,F_6)$ for the six declared invariants. The factorization $C=h_0(F_1,\ldots,F_6)$ is assumed without justification. By Theorem~\ref{thm:factorization}, this is equivalent to the unverified assertion that $C$ is constant on the level sets of $(F_1,\ldots,F_6)$.
\end{nonclaim}

\begin{nonclaim}[H4 verification status: failed]
No Lipschitz constant $L_h$ has been computed for any QICN decision rule. No decision margin $\Delta^{*}$ has been evaluated. The robustness condition $\Delta^{*}>L_h\sum_i\varepsilon_i$ has not been checked. Without H4, there is no guarantee that the adjudicator's verdict is stable under estimation error.
\end{nonclaim}

\subsection{Operational Estimator Formalization for Synthetic Fixtures}
\label{sec:operational-estimators}

The six QICN invariants declared in the synthetic fixture can be given a finite operational typing discipline. This is not a topological proof for general systems; it is an executable certificate format that makes the H2 and H4 quantities auditable inside a synthetic diagnostic manifest.

\begin{definition}[Operational latent invariant]\label{def:operational-invariant}
For a synthetic fixture, an \emph{operational latent invariant} is a map
\[
F_i^{\mathrm{op}}\colon\mathcal{P}\to Z_i
\]
where $\mathcal{P}$ is the finite set of admissible measurement points and $(Z_i,d_i)$ is a declared metric target interval or ray. The operational domain is the finite fixture point set, not the topological inverse limit $\mathcal{I}$ or a verified latent state space $X$.
\end{definition}

\begin{table}[h]
\centering
\caption{Operational invariant typing for the synthetic fixture family}
\label{tab:operational-invariants}
\begin{tabular}{llllc}
\hline
Invariant & Domain & Codomain $(Z_i,d_i)$ & $K_i^{\mathrm{op}}$ & $\varepsilon_i$ \\
\hline
identity\_channel\_lock & fixture points & $[0,1]$, $|a-b|$ & $1.0$ & $0.05$ \\
history\_alignment & fixture points & $[-1,1]$, $|a-b|$ & $2.0$ & $0.05$ \\
response\_phase & fixture points & $[0,2\pi)$, circular metric & $1.0$ & $0.05$ \\
gauge\_stability & fixture points & $[0,\infty)$, $|a-b|$ & $0.5$ & $0.05$ \\
intervention\_fidelity & fixture points & $[0,1]$, $|a-b|$ & $1.0$ & $0.05$ \\
factorization\_gap & fixture points & $[0,\infty)$, $|a-b|$ & $0.5$ & $0.05$ \\
\hline
\end{tabular}
\end{table}

\begin{nonclaim}[Declared operational constants are not derived Lipschitz constants]\label{nonclaim:declared-op-constants}
The values $K_i^{\mathrm{op}}$ in Table~\ref{tab:operational-invariants} are declared finite-fixture certificate bounds. They are not derived Lipschitz constants for latent invariants on an externally specified metric space $(X,d_X)$, and they do not verify the topological hypothesis H2 of Theorem~\ref{thm:bridge}. Their interpretation depends on the operational fixture metric and cannot be transferred to an unspecified latent space.
\end{nonclaim}

\begin{proposition}[Operational H2 gate]\label{prop:op-H2}
If each invariant in Table~\ref{tab:operational-invariants} is declared with $K_i^{\mathrm{op}}>0$, $\varepsilon_i>0$, and $\omega_i^{\mathrm{op}}\leq 2\varepsilon_i$, then the finite fixture satisfies the operational H2 certificate format. This does not imply that the true topological fiber bound
\[
\operatorname{diam}_X(\pi^{-1}(y)\cap A)\leq 2\varepsilon_i/K_i
\]
has been proved for any external latent state space $X$.
\end{proposition}

\begin{proof}
The executable gate checks the declared positivity of $K_i^{\mathrm{op}}$ and $\varepsilon_i$, and checks the finite certificate inequality $\omega_i^{\mathrm{op}}\leq 2\varepsilon_i$ for all six declared invariants. These are syntactic and arithmetic checks on the manifest certificate. They are sufficient to reject missing or inconsistent H2 certificates, but they do not construct $X$, $\pi$, $A$, or true topological constants $K_i$.
\end{proof}

\begin{nonclaim}[Finite sample checks do not verify universal fiber bounds]\label{nonclaim:finite-sample-fibers}
The topological H2 condition is quantified over every $y\in\pi(A)$ and involves the supremal oscillation of $F_i$ on the full fiber $\pi^{-1}(y)\cap A$. When $\pi(A)$ is infinite, this is a universal geometric condition, not an $n=8$ fixture check. The operational gate verifies only the finite manifest certificate; it does not prove $\omega_i(y)\leq 2\varepsilon_i$ for all latent observations in an external system.
\end{nonclaim}

\begin{proposition}[Operational H4 gate]\label{prop:op-H4}
For the finite synthetic adjudication manifest, define the pointwise support margin
\[
m(p)=|p.\mathrm{observed}_{\delta}-p.\mathrm{rival}|-|p.\mathrm{observed}_{\delta}-p.\mathrm{qicn}|.
\]
The executable operational margin is
\[
\Delta^{*,\mathrm{op}}=\min_{p\in\mathcal{P},\,m(p)>0}m(p),
\]
and the operational support function is assigned $L_h^{\mathrm{op}}=2$ because it is a difference of two absolute-value distance terms. The finite H4 gate passes only if
\[
\Delta^{*,\mathrm{op}}>L_h^{\mathrm{op}}\sum_i\varepsilon_i.
\]
\end{proposition}

\begin{proof}
Each absolute-value distance term is $1$-Lipschitz in its scalar prediction argument. The difference of two such terms is therefore bounded by Lipschitz constant $2$ for this finite scalar surrogate. The gate then performs the arithmetic comparison above. On the current fixture family, the H4 gate is allowed to fail; failure is evidence that the declared estimator-error budget is too large for the finite margin, not evidence for or against any external QICN claim.
\end{proof}

\begin{nonclaim}[Operational versus topological closure]
Definition~\ref{def:operational-invariant} and Propositions~\ref{prop:op-H2}--\ref{prop:op-H4} close only an executable certificate gap for synthetic fixtures. They do not prove H1, do not prove H3, do not show $C\in\sigma(F_1,\ldots,F_6)$, and do not establish that $F_i^{\mathrm{op}}$ are restrictions of genuine topological invariants $F_i\colon X\to Z_i$. The bridge theorem remains conditional for external systems.
\end{nonclaim}

\begin{nonclaim}[Fisher Information invalidation status]
Theorem~\ref{thm:fisher-asymmetry} is a conditional result about the mismatch between iid and residual-dependent Fisher Information geometry under its stated assumptions. It does not establish external support, does not validate QICN, and does not by itself adjudicate the v27 fixture outside the statistical diagnostics explicitly listed below.
\end{nonclaim}

\begin{nonclaim}[Prediction Jacobian is empirical, not derived]
The characterization of $J_{\mathrm{QICN}}$ as concentrated in the low-frequency eigenspace of $\bm{\Sigma}_u$ (and the resulting claim that the iid model ``actively mis\-leads'') is an empirical diagnostic on the v27 fixture, not a derived property of QICN. The v30 bridge theorem does not require this characterization and does not imply it. Future fixtures with different residual structure may exhibit different Jacobian alignments.
\end{nonclaim}

\begin{nonclaim}[AICc gain reversal is fixture-specific]
The reported AICc values (iid gain $+87.59$, Prais--Winsten gain $-48.59$, GLS gain $-48.68$) are diagnostics on the v27 fixture under the AR(1) model documented in N2. They are not derived from the v30 bridge theorem and do not constitute evidence about QICN's external validity. The sign reversal between iid and corrected estimators is consistent with the directional Jacobian bias documented in N1, but neither sign is interpreted here as support for or against QICN's external claims.
\end{nonclaim}

\begin{nonclaim}[AICc selection is conditional]
Proposici\'on~3.1 establishes that AICc selects between models only when the candidate set is exhaustive and the model class is correctly specified. Neither condition is verified for the QICN vs rival comparison on the v27 fixture. The reported AICc values must be read as descriptive diagnostics, not as model selection evidence.
\end{nonclaim}

\begin{corollary}[QICN bridge status: conjectural]
Since H2, H3, and H4 have not been verified for QICN, the conclusion of Theorem~\ref{thm:bridge} does not apply. The bridge remains unbuilt. The theorem provides four precise falsification conditions: (1) exhibit the topology of $X$, (2) compute $\omega_i(y)$ and verify the fiber diameter bound, (3) establish or refute $C\in\sigma(F_1,\ldots,F_k)$, (4) compute $L_h$ and verify $\Delta^{*}>L_h\sum_i\varepsilon_i$. Until all four are satisfied, no finite runner output constitutes mathematical evidence for any QICN claim.
\end{corollary}

\begin{corollary}[Statistical falsification: dependence and leakage gates]
Independent of the topological conditions H1--H4, the statistical validity of the adjudicator requires the residual process to satisfy the iid or GLS likelihood framework of Section~\ref{sec:gls}. The v27 fixture fails this requirement catastrophically ($\mathrm{DW}\approx0.038$) and is also blocked by structural leakage and rival-adequacy gates. Older hardcoded claims of a specific AR(1) sign reversal from $+87.59$ to $-50.38$ are not admissible unless reproduced under the current centered-rho executable runner. Any fixture that passes only under iid assumptions while failing the dependence, leakage, or rival gates is statistically void.
\end{corollary}

\begin{nonclaim}[No hidden global bridge]
Theorem~\ref{thm:bridge} is not a derivation of $M_{\Omega}=+\infty$, global separator atomicity, external support, consciousness, phenomenality, identity transfer, or bridge-burden closure from RSS, AIC, Gaussian likelihood, or a finite fixture. It is a conditional theorem with four explicit, independently falsifiable hypotheses. None have been verified for QICN. The theorem stands; the bridge does not.
\end{nonclaim}

%% ============================================================
\section{Exact GLS Implementation and Computational Verification}
\label{sec:gls-implementation}
%% ============================================================

This section documents the v30 GLS integration, providing the exact computational formulas that underlie the statistical results of Section~\ref{sec:gls}.

\begin{definition}[AR(1) profile likelihoods used by the runner]\label{def:gls-nll}
For a residual vector $\bm{u}$ under AR(1) with parameter $\rho$, v30 distinguishes two finite-record profile likelihoods.

\textbf{Prais--Winsten innovation form.} Let
\[
v_1=\sqrt{1-\rho^2}\,u_1,\qquad v_t=u_t-\rho u_{t-1}\quad(t\geq2),
\]
and let $\hat{\sigma}_\varepsilon^2=\bm{v}^T\bm{v}/n$. The innovation negative log-likelihood used by \texttt{correctedGaussian\allowbreak{}Information} is
\begin{equation}\label{eq:pw-nll}
\mathrm{NLL}_{\mathrm{PW}}\;=\;\frac{n}{2}\bigl(\ln(2\pi\hat{\sigma}_\varepsilon^2)+1\bigr)-\frac{1}{2}\ln(1-\rho^2).
\end{equation}
The final term is the Prais--Winsten first-observation Jacobian correction.

\textbf{Full determinant GLS form.} Let $\bm{V}$ be the AR(1) correlation matrix with entries $V_{st}=\rho^{|s-t|}$ and let $\hat{\sigma}^2=\bm{u}^T\bm{V}^{-1}\bm{u}/n$. The exact GLS profile negative log-likelihood used by \texttt{glsGaussianInformation} is
\begin{equation}\label{eq:gls-nll-det}
\mathrm{NLL}_{\mathrm{GLS}}\;=\;\frac{n}{2}\bigl(\ln(2\pi\hat{\sigma}^2)+1\bigr)+\frac{1}{2}\ln|\bm{V}|
\;=\;\frac{n}{2}\bigl(\ln(2\pi\hat{\sigma}^2)+1\bigr)+\frac{n-1}{2}\ln(1-\rho^2).
\end{equation}
\end{definition}

\begin{proposition}[Scaled AR(1) inverse and quadratic form]\label{prop:quadratic-form}
For the AR(1) correlation matrix $\bm{V}_{st}=\rho^{|s-t|}$,
\begin{equation}\label{eq:V-inv}
\bm{V}^{-1}\;=\;\frac{1}{1-\rho^2}\begin{pmatrix}
1 & -\rho & 0 & \cdots & 0 & 0 \\
-\rho & 1+\rho^2 & -\rho & \cdots & 0 & 0 \\
0 & -\rho & 1+\rho^2 & \cdots & 0 & 0 \\
\vdots & \vdots & \vdots & \ddots & \vdots & \vdots \\
0 & 0 & 0 & \cdots & 1+\rho^2 & -\rho \\
0 & 0 & 0 & \cdots & -\rho & 1
\end{pmatrix}.
\end{equation}
Consequently,
\begin{equation}\label{eq:pw-quadratic}
\bm{u}^T\bm{V}^{-1}\bm{u}\;=\;\frac{(1-\rho^2)u_1^2+\sum_{t=2}^n(u_t-\rho\,u_{t-1})^2}{1-\rho^2}.
\end{equation}
\end{proposition}

\begin{proof}
The displayed tridiagonal matrix is the standard inverse of the AR(1) correlation matrix after scaling by $(1-\rho^2)^{-1}$. Multiplying $\bm{u}$ by the unscaled tridiagonal part gives the numerator
\[
(1-\rho^2)u_1^2+\sum_{t=2}^n(u_t-\rho u_{t-1})^2.
\]
The factor $(1-\rho^2)^{-1}$ belongs to $\bm{V}^{-1}$ and must not be omitted. This is the scaling implemented by \texttt{ar1QuadraticForm}.
\end{proof}

\begin{proposition}[Determinant correction]\label{prop:determinant}
The determinant of the AR(1) correlation matrix is
\begin{equation}\label{eq:log-det}
|\bm{V}|=(1-\rho^2)^{n-1},\qquad
\ln|\bm{V}|=(n-1)\ln(1-\rho^2).
\end{equation}
Thus the full GLS profile likelihood contains the determinant term $\frac{n-1}{2}\ln(1-\rho^2)$.
\end{proposition}

\begin{proof}
This is the classical determinant formula for the Toeplitz AR(1) correlation matrix. Equivalently, the determinant follows from the LU factorization of the scaled tridiagonal inverse in \eqref{eq:V-inv}. No residual determinant term is absorbed into $\hat{\sigma}^2$ after profile substitution; the determinant term remains explicit in \eqref{eq:gls-nll-det}.
\end{proof}

\begin{corollary}[PW versus GLS diagnostic paths]\label{cor:gls-co-diff}
The Prais--Winsten innovation path and the full determinant GLS path are both valid internal diagnostics only under the declared AR(1) model. They use different scale conventions and should be reported separately. For the current executable fixture, both corrected paths produce negative AICc gain, so the iid gain remains statistically void.
\end{corollary}

\begin{nonclaim}[v30 computational verification]
The v30 adjudicator reports iid AICc, Prais--Winsten corrected AICc, and exact GLS AICc. Current executable values for the v27 fixture are approximately: iid gain $+87.59$, Prais--Winsten corrected gain $-48.59$ before the v32 Jacobian correction, and GLS gain $-48.68$. These are internal synthetic diagnostics only and do not certify external support or bridge closure.
\end{nonclaim}

\end{document}
